\section{Discusión y amenazas a la validez (Entrega 4)}
\label{sec:discusion-e4}

\subsection{Interpretación de los resultados}

El sistema cumple de forma sólida las características de calidad relacionadas con
\textbf{estructura y verificabilidad del código} (arquitectura hexagonal con dependencia del
dominio hacia adentro, seis patrones GoF con problema real documentado, complejidad ciclomática
media de 1.8, capa unitaria de la pirámide de pruebas con 438 pruebas automatizadas ejecutadas en
el último ciclo) y las relacionadas con \textbf{compatibilidad de cliente} (la misma suite E2E pasa en tres
motores de navegador distintos, y el código respeta su API mínima declarada). Donde el sistema
\emph{no} alcanza el umbral objetivo es, de forma consistente, bajo \textbf{concurrencia real}:
tanto la eficiencia de desempeño como la fiabilidad se degradan por una limitación de la edición
gratuita del motor de base de datos (Sección~\ref{sec:iso25010}), no por un defecto de diseño de
la aplicación. Esta distinción —entre ``el diseño no soporta esta carga'' y ``el entorno de
demostración no soporta esta carga''— es central para interpretar correctamente los resultados: la
arquitectura (Raft con quorum de mayoría, \emph{stateless} en cada microservicio, saga por
coreografía sin coordinador central) sí está diseñada para escalar horizontalmente; lo que falta
es la licencia o el plan de CockroachDB que levante el límite artificial de transacciones
concurrentes.

\subsection{Alcance y limitaciones}

Este proyecto es un sistema de demostración académica, no un despliegue de producción: corre en
una sola máquina de desarrollo con Docker Desktop, sobre datos sintéticos, sin balanceo de carga
real entre réplicas, y sin el hardware dedicado que un ISP real destinaría a este tipo de
sistema. Las conclusiones sobre desempeño y fiabilidad son válidas \emph{para este entorno}, y se
documentan explícitamente como tales en vez de generalizarse sin matización a cualquier
despliegue posible.

\subsection{Amenazas a la validez del proyecto (a nivel general)}

Más allá de las amenazas específicas del experimento de carga (ya discutidas en
Sección~\ref{sec:iso25010}), a nivel de todo el proyecto:

\textbf{Internas}:
\begin{enumerate}[nosep]
  \item \textbf{Un solo desarrollador principal por capa}: cada integrante del equipo se
        concentró en un área (backend, observabilidad, contratos, documentación), lo que agiliza
        el desarrollo pero concentra el conocimiento profundo de cada módulo en una sola persona
        —mitigado parcialmente por la documentación exhaustiva (ADRs, esta guía) que permite a
        cualquier integrante razonar sobre módulos que no escribió directamente.
  \item \textbf{Verificación local, no en el entorno de ejecución final}: buena parte de la
        verificación de este ciclo (pruebas de backend con Testcontainers, cobertura, complejidad
        ciclomática) se corrió en un contenedor Maven anidado dentro de Docker Desktop en
        Windows, un entorno con una limitación conocida de redes anidadas (\emph{Docker outside
        of Docker}) que impidió ejecutar una prueba de integración específica localmente —mitigado
        porque esa misma prueba sí corre en verde en el \emph{runner} de GitHub Actions, que
        ejecuta Maven directo sobre la máquina virtual sin anidar Docker.
  \item \textbf{Datos de demostración generados sintéticamente}: los 504+ tickets de prueba se
        generaron y rebalancearon con un script SQL propio, no provienen de un ISP real —mitigado
        documentando explícitamente el script de generación (\texttt{resultados/rebalance\_demo\_tickets.sql})
        para que la distribución de estados sea auditable y reproducible.
\end{enumerate}

\textbf{Externas}:
\begin{enumerate}[nosep]
  \item \textbf{Generalización a otros ISP}: el dominio (zonas de cobertura, categorías de
        incidencia, políticas de SLA) se modeló según supuestos razonables del equipo, no según
        un proceso de negocio documentado de un operador real —una limitación externa reconocida
        desde la Entrega 1.
  \item \textbf{Generalización a escala de producción}: las cifras de esta evaluación (latencia,
        tasa de error, cobertura) corresponden a un solo nodo de desarrollo con datos sintéticos;
        no se puede extrapolar directamente el comportamiento a un despliegue con tráfico real de
        miles de usuarios concurrentes sin repetir el experimento en ese entorno.
\end{enumerate}

\subsection{Reflexión ética}

El sistema procesa datos personales reales de abonados de un ISP: ubicación de la incidencia,
historial de contacto, y las credenciales de acceso de clientes, técnicos y administradores. El
Código de Ética y Conducta Profesional de la ACM~\cite{acm2018ethics} establece, entre sus
principios generales, evitar daño y respetar la privacidad; frente a esos principios, el equipo
identifica lo siguiente:

\textbf{Ya aplicado en el sistema}: las contraseñas nunca se almacenan en texto plano
(\texttt{BCryptPasswordEncoder}, factor de costo por defecto de Spring Security); los tokens de
sesión (JWT) tienen expiración corta y un mecanismo de revocación vía \emph{refresh tokens}; y
ningún \emph{endpoint} de administración de la plataforma (\texttt{actuator}) expone información
sensible del entorno de ejecución (variables, volcado de memoria).

\textbf{Reconocido como pendiente, no resuelto}: la revisión de seguridad de esta entrega
(Sección~\ref{sec:iso25010}) encontró que el sistema no protege el endpoint de inicio de sesión
contra intentos repetidos de adivinar una contraseña. Frente a datos reales de personas —técnicos
de campo cuya ubicación queda registrada, clientes cuyo domicilio se infiere de su zona de
cobertura— esta es una responsabilidad ética concreta, no solo un requisito técnico: el equipo
documenta esta brecha de forma explícita, en vez de omitirla, precisamente porque una evaluación
honesta de los riesgos reales del sistema es un prerrequisito para poder mitigarlos antes de un
eventual uso con usuarios reales.
